% Standalone mathematical body; standard amsmath/amsthm notation.
% Suggested macros: \F=\mathbb F_q, \Pbb=\mathbb P, \E=\mathbb E,
% \Alt=\operatorname{Alt}, \TV=d_{\rm TV}.

\section{A hidden symplectic form: the statistical experiment}

Fix a prime power $q$, an even integer $n\ge4$, and $V=\mathbb F_q^n$.
An alternating form means a bilinear form $H$ satisfying $H(v,v)=0$ for
every $v$, including in characteristic two. Put
\begin{gather*}
 D=\binom n2,\qquad d=D-1,\qquad
 \gamma_q=\prod_{j\ge1}(1-q^{-j}),\\
 c_n=\prod_{j=2}^{n/2}(1-q^{-(2j-1)}),\qquad
 c_q=\prod_{j\ge2}(1-q^{-(2j-1)}).
\end{gather*}
Empty products are one. Let $\mathcal H_n$ be the set of nondegenerate
alternating forms modulo multiplication by a nonzero scalar. Its cardinality is
\begin{equation}\label{hidden:eq:M}
 M_n=q^d c_n.
\end{equation}
Let $\mathcal G_n=\operatorname{Gr}_2(V)$ be the set of two-dimensional
subspaces. A member $K$ is $H$-isotropic if $H|_{K\times K}=0$.
For every nondegenerate $H$, the fraction of $\mathcal G_n$ that is
$H$-isotropic is
\begin{equation}\label{hidden:eq:p}
 p_n=\frac{q^{n-2}-1}{q^{n-1}-1}<q^{-1}.
\end{equation}

We compare two laws on ordered $t$-tuples of two-planes.
Under $Q_{n,t}$ all the planes are independent and uniform on
$\mathcal G_n$. Under $P_{n,t}$, first choose $[H]$ uniformly from
$\mathcal H_n$ and then, conditional on $H$, choose the planes independently
and uniformly among its isotropic two-planes. The hidden form is chosen
\emph{once} and shared by all the observations. Different scalar multiples
of $H$ are deliberately identified because they produce the same experiment.

For observed planes $K_1,\ldots,K_t$, define the linear solution space
\begin{equation}\label{hidden:eq:solution}
 \mathcal N(K_1,\ldots,K_t)
 =\{G\in\Alt(V):G|_{K_i\times K_i}=0\text{ for every }i\}
\end{equation}
and let $Z$ be the number of projective nondegenerate forms in
$\mathcal N$. Each plane contributes one homogeneous linear equation in the
$D$ entries above the diagonal of $G$. Thus $\mathcal N$ is obtained by
ordinary Gaussian elimination. Under $P_{n,t}$ it contains the line
$\langle H\rangle$.

\begin{theorem}[A sharp additive sample threshold]\label{hidden:thm:threshold}
For every fixed $q$, as even $n\to\infty$, the experiment has threshold
\[
 t=\binom n2-1.
\]
More precisely, for any positive real sequence $w_n\to\infty$:
\begin{enumerate}
\item If $0\le t_n\le d-w_n$, then
\[
 d_{\rm TV}(P_{n,t_n},Q_{n,t_n})\longrightarrow0.
\]
The optimal probability of recovering the hidden projective form from
these observations also tends to zero.
\item If $t_n\ge d+w_n$, then
\[
 d_{\rm TV}(P_{n,t_n},Q_{n,t_n})\longrightarrow1.
\]
Uniformly in the nondegenerate hidden form $H$, with probability tending
to one under its conditional observation law,
\[
 \mathcal N(K_1,\ldots,K_{t_n})=\langle H\rangle.
\]
Thus Gaussian elimination recovers the hidden form up to scalar above
the same threshold.
\end{enumerate}
All statements hold in every characteristic. They require no hypothesis
about zeros of an $L$-function.
\end{theorem}

The theorem concerns an additive diverging gap, not merely a fixed
relative gap $\varepsilon D$. It does not claim a limiting distribution
for $Z$ throughout the bounded critical window.

\begin{theorem}[The bounded window and its second moment]\label{hidden:thm:window}
For each fixed integer $a$, at $t=d+a$,
\begin{equation}\label{hidden:eq:chi-limit}
 \chi^2(P_{n,d+a}\Vert Q_{n,d+a})\longrightarrow c_q^{-1}q^a.
\end{equation}
Consequently, for each fixed nonnegative integer $w$,
\begin{align}
 \limsup_{n\to\infty}d_{\rm TV}(P_{n,d-w},Q_{n,d-w})
 &\le\tfrac12c_q^{-1/2}q^{-w/2},\label{hidden:eq:below-TV}\\
 \limsup_{n\to\infty}\Pr(\text{optimal recovery at }d-w)
 &\le c_q^{-1}q^{-w},\label{hidden:eq:below-recovery}\\
 \liminf_{n\to\infty}d_{\rm TV}(P_{n,d+w},Q_{n,d+w})
 &\ge 1-c_q q^{-w},\label{hidden:eq:above-TV}\\
 \limsup_{n\to\infty}\sup_H
 \Pr_H\{\mathcal N\ne\langle H\rangle\text{ at }d+w\}
 &\le\frac{q^{-w}}{q-1}.\label{hidden:eq:above-recovery}
\end{align}
Bounds with a right side outside $[0,1]$ can of course be clipped to that
interval. In particular the window needed for any fixed accuracy is
bounded independently of $n$.
\end{theorem}

\section{Exact likelihood and finite-dimensional estimates}

Set
\[
 \rho_q=\frac{2q-1}{q^2}<1,\qquad
 K_q=4q(q-1),\qquad
 A_q=\frac{q}{\gamma_q c_q}.
\]
The constants are convenient rather than optimized.

\begin{proposition}[Finite bounds]\label{hidden:prop:finite}
For every $n,t$ as above and every integer $0\le R<n$,
\begin{align}
 \chi^2(P_{n,t}\Vert Q_{n,t})
 &\le e^{K_qtq^{-R}}-1+\frac1{M_np_n^t}
       +A_q q^{-\binom{n-R}{2}}\left(\frac{\rho_q}{p_n}\right)^t,
       \label{hidden:eq:finite-chi}\\
 Q_{n,t}(Z>0)&\le M_np_n^t,\label{hidden:eq:null-union}\\
 \Pr(\text{optimal recovery})&=\frac{Q_{n,t}(Z>0)}{M_np_n^t}
                         \le\frac1{M_np_n^t},
                         \label{hidden:eq:Bayes}\\
 \sup_H\Pr_H\{\mathcal N\ne\langle H\rangle\}
 &\le \frac{q^d}{q-1}\,[p_ne^{K_qq^{-R}}]^t
  +\frac{q^{D-\binom{n-R}{2}}}{(q-1)\gamma_q}\rho_q^t.
  \label{hidden:eq:finite-recovery}
\end{align}
The likelihood itself is
\begin{equation}\label{hidden:eq:likelihood}
 \frac{dP_{n,t}}{dQ_{n,t}}=\frac{Z}{M_np_n^t}.
\end{equation}
In particular $P_{n,1}=Q_{n,1}$ exactly: one observation has no information
for testing the mixture against the null.
\end{proposition}

\subsection{Counting forms and isotropic planes}

The general linear group acts transitively on the nondegenerate alternating
forms, with stabilizer the symplectic group. The standard basis counts
give
\[
 |\GL_n(q)|=q^{n^2}\prod_{j=1}^n(1-q^{-j}),\qquad
 |\Sp_n(q)|=q^{n(n+1)/2}\prod_{j=1}^{n/2}(1-q^{-2j}).
\]
Their quotient is
$q^D\prod_{j=1}^{n/2}(1-q^{-(2j-1)})$.
Division by $q-1$ yields \eqref{hidden:eq:M}. These formulas can also be
obtained by choosing successive hyperbolic pairs, so they use no division
by two.

To count isotropic planes, choose an ordered independent pair $(u,v)$.
For each nonzero $u$, the functional $H(u,\cdot)$ is nonzero and its
kernel has $q^{n-1}$ elements. Removing the $q$ elements of $\langle u\rangle$
leaves $q^{n-1}-q$ choices of $v$, out of $q^n-q$ unrestricted choices.
This gives \eqref{hidden:eq:p}.

For a given tuple each compatible $[H]$ has conditional likelihood
$p_n^{-t}$ relative to $Q_{n,t}$. Averaging proves
\eqref{hidden:eq:likelihood}; also $\mathbb E_Q Z=M_np_n^t$, which proves
\eqref{hidden:eq:null-union}. The posterior under the planted model is
uniform among the $Z$ compatible forms. Therefore the optimal Bayes
success is
\[
 \mathbb E_P\frac1Z
 =\sum_{\{Z>0\}}Q(\boldsymbol K)\frac{Z}{M_np_n^t}\frac1Z
 =\frac{Q(Z>0)}{M_np_n^t},
\]
proving \eqref{hidden:eq:Bayes}. Transitivity of $\GL(V)$ on the
two-planes shows that a single plane is uniform after averaging $H$;
equivalently, the number of compatible forms is constant on $\mathcal G_n$.

\subsection{Two-form incidence}

\begin{lemma}[Common isotropic planes]\label{hidden:lem:pair}
Fix a nondegenerate alternating $H$ and any alternating $G$. Put
\[
 e(H,G)=\sum_{a\in\mathbb F_q}
    \bigl(q^{\dim\ker(G-aH)}-1\bigr),\qquad
 r(H,G)=\frac{e(H,G)}{q^n-1}.
\]
The probability that a uniform two-plane is isotropic for both forms is
\begin{equation}\label{hidden:eq:pair}
 p_2(H,G)=
 \frac{e(H,G)(q^{n-1}-q)+(q^n-1-e(H,G))(q^{n-2}-q)}
 {(q^n-1)(q^n-q)}.
\end{equation}
If $G$ is not proportional to $H$, then
\begin{equation}\label{hidden:eq:rho}
 \Pr_H\{K\text{ is }G\text{-isotropic}\}
       =\frac{p_2(H,G)}{p_n}\le\rho_q.
\end{equation}
If $s=\min_{a\in\mathbb F_q}\operatorname{rank}(G-aH)$, then
\begin{equation}\label{hidden:eq:far}
 \frac{p_2(H,G)}{p_n^2}\le1+K_q q^{-s}.
\end{equation}
\end{lemma}

\begin{proof}
For each nonzero $u$, the linear forms $H(u,\cdot)$ and $G(u,\cdot)$
are dependent exactly when $(G-aH)u=0$ for one scalar $a$.
The scalar is unique because $H$ is nondegenerate, so the number of these
$u$ is precisely $e(H,G)$. Their common kernel has dimension $n-1$;
otherwise it has dimension $n-2$. Both kernels contain $u$ because both
forms are alternating. Counting eligible $v$ gives \eqref{hidden:eq:pair}.

If $G$ is not proportional, every $G-aH$ is a nonzero alternating form,
so has even rank at least two. Thus
\[
 r(H,G)\le\frac{q(q^{n-2}-1)}{q^n-1}<q^{-1}.
\]
Conditional on $H$-isotropy, the probability of $G$-isotropy is
\[
 r(H,G)+(1-r(H,G))\frac{q^{n-2}-q}{q^{n-1}-q}
 \le\frac1q+\left(1-\frac1q\right)\frac1q=\rho_q,
\]
which proves \eqref{hidden:eq:rho}.

For \eqref{hidden:eq:far}, put $x=q^{n-1}$. Direct substitution gives
\begin{equation}\label{hidden:eq:C-explicit}
 \frac{p_2(H,G)}{p_n^2}
 =1+\frac{(q-1)x[(x-1)r(H,G)-(q-1)]}{(x-q)^2}.
\end{equation}
Since $n\ge4$, $x(x-1)/(x-q)^2\le2$; for example
$x\ge q^3$ implies $x/(x-q)\le(1-q^{-2})^{-1}\le4/3$.
Also
$r(H,G)\le q q^{n-s}/(q^n-1)\le2q q^{-s}$.
Dropping the negative term in \eqref{hidden:eq:C-explicit} proves
\eqref{hidden:eq:far}.
\end{proof}

The pair-incidence mechanism is classical finite geometry; it is closely
related to the eigenvector counts used for symplectic Grassmann codes by
Cardinali and Giuzzi (2016). It is included to make all statistical
estimates independently checkable. The new claim here concerns the
shared hidden-form experiment and its sharp sample threshold.

\subsection{Few forms are close to a scalar multiple}

\begin{lemma}[A rank-variety bound]\label{hidden:lem:rank}
The number $B_{n,R}$ of alternating forms on $V$ having rank at most $R$
satisfies
\begin{equation}\label{hidden:eq:rankcount}
 B_{n,R}\le\gamma_q^{-1}q^{D-\binom{n-R}{2}}.
\end{equation}
For a fixed nondegenerate $H$, a uniform nondegenerate $G$ has
\[
 \Pr\{\min_a\operatorname{rank}(G-aH)\le R\}
 \le A_q q^{-\binom{n-R}{2}}.
\]
\end{lemma}

\begin{proof}
Every rank-at-most-$R$ form has a kernel containing some $(n-R)$-space
$W$. There are at most
$\gamma_q^{-1}q^{R(n-R)}$ such $W$: use the product formula for a Gaussian
binomial coefficient, bound its numerator by one and its denominator
below by $\gamma_q$. Forms vanishing on $W$ descend to alternating forms
on the $R$-dimensional quotient, of which there are $q^{R(R-1)/2}$.
This proves \eqref{hidden:eq:rankcount}, with harmless overcounting.
For nondegenerate $G$ and $R<n$, the scalar $a=0$ cannot contribute.
The union over the remaining $q-1$ scalars has at most $(q-1)B_{n,R}$
elements. The number of nondegenerate $G$ is $(q-1)M_n$, so division
gives $B_{n,R}/M_n\le q(\gamma_qc_q)^{-1}q^{-\binom{n-R}{2}}$.
\end{proof}

\subsection{Proof of the chi-square estimate}

Squaring \eqref{hidden:eq:likelihood} and averaging under $Q$ gives the
exact identity
\begin{equation}\label{hidden:eq:secondmoment}
 1+\chi^2(P_{n,t}\Vert Q_{n,t})
 =\mathbb E_{[H],[G]\in\mathcal H_n}
       \left(\frac{p_2(H,G)}{p_n^2}\right)^t.
\end{equation}
The expectation can equivalently be over all nondegenerate forms,
because every projective class has the same size. Equal projective
classes have probability $1/M_n$ and contribute exactly
$1/(M_np_n^t)$. For all other pairs with
$\min_a\operatorname{rank}(G-aH)>R$, use \eqref{hidden:eq:far} and
$(1+z)^t\le e^{tz}$. On the remaining nonproportional pairs use
\eqref{hidden:eq:rho} and \cref{hidden:lem:rank}. Subtracting one yields
\eqref{hidden:eq:finite-chi}.

\subsection{Proof of the recovery estimate}

Fix $H$. A competitor $G\notin\langle H\rangle$ gives the same event
on $H$-isotropic planes as $aG+bH$, for every $a\ne0$ and every $b$.
Thus distinct events can be indexed by the projective points of the
$d$-dimensional quotient $\Alt(V)/\langle H\rangle$. There are
$(q^d-1)/(q-1)$ such points, and each represents exactly $q(q-1)$
forms $aG+bH$.

The quantity $\min_b\operatorname{rank}(G-bH)$ is invariant under these
changes of representative. At most $qB_{n,R}$ representatives have
this minimum at most $R$, by the union over $q$ translates of the
rank-at-most-$R$ set. Consequently at most $B_{n,R}/(q-1)$ quotient
projective points are near. For each of them the probability of surviving
all $t$ observations is at most $\rho_q^t$. For the far points the
survival probability is at most $[p_ne^{K_qq^{-R}}]^t$ by
\eqref{hidden:eq:far}. The union bound and
\eqref{hidden:eq:rankcount} prove \eqref{hidden:eq:finite-recovery}.
If no competitor survives, the solution space is exactly
$\langle H\rangle$, which makes recovery by Gaussian elimination valid.

\section{From the finite bounds to the sharp threshold}

Take $R=\lfloor n/5\rfloor$. Then
$tq^{-R}\to0$ whenever $t=O(n^2)$. Put
\[
 \lambda_q=\log_q(q\rho_q)=\log_q(2-q^{-1}).
\]
For $q=2$, $\lambda_2=\log_2(3/2)<16/25$; for $q\ge3$,
$\lambda_q<\log_3 2<16/25$. Also
\[
 \binom{n-R}{2}=\left(\frac{16}{25}+o(1)\right)D,
 \qquad p_n=q^{-1}(1+O_q(q^{2-n})),\qquad
 M_n=q^d(c_q+o(1)).
\]
Thus the near term in \eqref{hidden:eq:finite-chi}, uniformly for
$0\le t\le d+O(1)$, is at most
\[
 q^{-\{16/25-\lambda_q+o(1)\}D}=o(1).
\]
The same calculation makes the near term in
\eqref{hidden:eq:finite-recovery} tend to zero for $t=d+O(1)$.

For $t=d+a$ with fixed $a$, the proportional term tends to
$c_q^{-1}q^a$. The finite chi-square bound therefore gives the upper
bound in \eqref{hidden:eq:chi-limit}. To obtain the matching lower bound,
\eqref{hidden:eq:C-explicit} with $r\ge0$ gives
\[
 \frac{p_2(H,G)}{p_n^2}\ge1-\delta_n,\qquad
 \delta_n=\frac{(q-1)^2q^{n-1}}{(q^{n-1}-q)^2}=O_q(q^{1-n}).
\]
The total contribution of the unequal projective pairs to
\eqref{hidden:eq:secondmoment} is consequently at least
$(1-M_n^{-1})(1-\delta_n)^t\to1$. Its upper bound tends to one by
the preceding near/far estimates. This proves the exact limit
\eqref{hidden:eq:chi-limit}.

Cauchy--Schwarz gives
$d_{\rm TV}(P,Q)\le\tfrac12\sqrt{\chi^2(P\Vert Q)}$.
Together with \eqref{hidden:eq:Bayes}, this gives the two below-threshold
bounds in \cref{hidden:thm:window}. The event $\{Z>0\}$ has planted
probability one, so
$d_{\rm TV}(P,Q)\ge1-Q(Z>0)\ge1-M_np_n^t$.
This proves \eqref{hidden:eq:above-TV}. Finally the far term of
\eqref{hidden:eq:finite-recovery} at $t=d+w$ tends to
$q^{-w}/(q-1)$, proving \eqref{hidden:eq:above-recovery}.

The assertions with a diverging additive gap follow either directly from
the finite bounds or by monotonicity. Additional observations cannot
decrease total variation or optimal recovery probability, and cannot
increase the probability that the linear solution space has dimension
greater than one. For example, for every fixed $w$ the experiment at
$t_n\le d-w_n$ is eventually a projection of the one at $d-w$; let
$n\to\infty$ and then $w\to\infty$. The other assertions are identical.

\section{The connection with symplectic compression}

The current ProveIt article's theorem \texttt{spn:thm:wide} gives the law
of a wide Lagrangian projection with a nondegenerate source form.
Specialize its source dimension to $n$, its target row dimension to
$n-2$, and its ambient symplectic dimension to $2(n-2)$. A uniform
symplectic transformation produces an $(n-2)\times n$ block $A$.
Condition on $A$ having full row rank. The exact fibre formula is constant
on the supported matrices of that rank and vanishes unless
$\ker A$ is $H$-isotropic. Therefore $\ker A$ is uniform among the
$H$-isotropic two-planes. For a uniform independent-entry block $U$,
conditioned on full row rank, $\ker U$ is uniform on all two-planes.

Given the kernel, both conditional block laws are uniform on the matrices
having that kernel. Hence passing from each block to its kernel loses no
information for either testing or recovery in these conditioned
experiments. Choosing the same unknown $[H]$ for the entire collection
of blocks gives exactly the experiment of \cref{hidden:thm:threshold}.
This supplies a sharp solution to one explicitly delimited case of the
article's question \texttt{spn:q:symplectic}: correlated observations
with an unknown shared source form.

For completeness, the needed fibre count has a short independent proof.
Put $r=n-2$ and use the ambient form
$\left(\begin{smallmatrix}0&I_r\\-I_r&0\end{smallmatrix}\right)$.
The symplectic group is transitive on the $2r\times n$ frames with Gram
matrix $H$: their nondegenerate images have isometric orthogonal
complements, so an isometry of their images extends to the ambient space.
Consequently the frame $(A;B)$ is uniform subject to
\[
 A^{\mathsf T}B-B^{\mathsf T}A=H.
\]
Every solution is automatically a frame because a vector in the kernel of
$(A;B)$ lies in the radical of $H$. For full-row-rank $A$, row and source
basis changes put $A$ in the form $(I_r\ 0)$. Write $B=(B_1\ B_2)$ and
split $H$ into blocks of sizes $r,2$. The equations become
\[
 B_1-B_1^{\mathsf T}=H_{11},\qquad B_2=H_{12},\qquad H_{22}=0.
\]
Thus the fibre is empty unless the kernel is isotropic, and otherwise has
exactly $q^{r(r+1)/2}$ members, from the diagonal and one triangle of
$B_1$. For each two-plane $K$, the number of full-row-rank $A$ with kernel
$K$ is $|\GL_r(q)|$, the number of isomorphisms $V/K\to\mathbb F_q^r$.
This proves all the uniformity assertions without importing a counting
claim from the source article. The full-row-rank probability is
independent of $H$, either by source congruence or by these counts, so
conditioning each observation leaves the uniform prior on $[H]$ unchanged.

The conditioning is part of the mathematical statement. No conclusion
about the sample threshold for unconditioned blocks is inferred here.
The generic group-counting and pair-incidence inputs are classical; the
contribution is the hidden-form threshold, its additive window, the
matching recovery statement, and their exact transfer to the conditioned
compression model.
